\chapter*{KẾT LUẬN VÀ HƯỚNG PHÁT TRIỂN}
\addcontentsline{toc}{chapter}{KẾT LUẬN VÀ HƯỚNG PHÁT TRIỂN}

\subsection{Kết luận}

Khóa luận xây dựng hai bộ dữ liệu đánh giá và đối chiếu hai kiến trúc gọi công cụ tiếng Việt trên Core Benchmark, CustomTools-VI và stress test. Kết quả cho thấy sự đánh đổi giữa độ chính xác trích xuất, khả năng xử lý schema mới, độ trễ ghi nhận và bộ nhớ sử dụng trong các điều kiện thực nghiệm đã mô tả.

Về dữ liệu, Canonical Core Benchmark gồm 77,028 cặp bản ghi song ngữ trên 4,421 công cụ; CustomTools-VI gồm 8,000 mẫu, trong đó 3,200 mẫu âm tính (40\%). Riêng mỗi tập kiểm thử \texttt{test_seen} và \texttt{test_unseen} cân bằng 50\% mẫu âm tính. Trên hai tập kiểm thử này, SLM E3 cho thấy thiên kiến kích hoạt công cụ quá mức, còn E4 đạt Non-FC Recall 100.0\%. Kết quả gợi ý dữ liệu bản địa có ích cho hành vi từ chối trong protocol hiện tại, song chưa xác định riêng đóng góp của từng thành phần huấn luyện.

Qwen3.5-4B E4 đạt ArgA 87.92\% trên CustomTools-VI seen và 86.75\% trên unseen, cao hơn GPT-5.6 Luna trong phép đánh giá này. Ở stress test, SLM gặp CUDA Out of Memory tại $N \ge 500$ trên GPU 16GB, còn Method 2 xử lý được $N=1000$ với ArgA 84.00\%, P50 107.68 ms và PyTorch allocated VRAM khoảng 3.21 GiB. Các số đo độ trễ của hai phương pháp được thu thập theo quy cách riêng (đơn truy vấn và batch), vì vậy không quy đổi thành hệ số tăng tốc tương đối.

Đóng góp của khóa luận là bộ dữ liệu đánh giá tiếng Việt có tập công cụ mới và mẫu không gọi công cụ, cùng phép đối chiếu hai kiến trúc trên độ chính xác, khả năng từ chối và độ bền theo quy mô danh mục. Các kết quả chỉ có giá trị trong bộ dữ liệu, phần cứng và giao thức đo đã báo cáo.

\subsection{Hướng phát triển trong tương lai}

Để kiểm tra khả năng khái quát của các kết quả, nghiên cứu tiếp theo cần lặp thực nghiệm qua nhiều seed, đo hai kiến trúc bằng cùng giao thức độ trễ và đánh giá chi phí trên các cấu hình triển khai thực tế.

Một hướng mở rộng là gọi công cụ trong hội thoại đa lượt. Hệ thống cần theo dõi trạng thái đối thoại, lưu các tham số đã được cung cấp và hỏi làm rõ khi thiếu thông tin trước khi tạo lời gọi. Hướng này cần bộ dữ liệu và độ đo riêng cho tính nhất quán của tham số qua nhiều lượt.

Kiến trúc lai cũng cần được kiểm chứng: Bi-Encoder chọn Top-$K$ công cụ, sau đó SLM sinh tham số trên tập ứng viên đã thu hẹp. Phép thử cần đo đồng thời độ chính xác, độ trễ và bộ nhớ theo cùng giao thức, đặc biệt trên công cụ chưa từng gặp và danh mục lớn.

Cuối cùng, nghiên cứu có thể mở rộng sang thực thi công cụ trong môi trường cô lập và xử lý phản hồi lỗi từ API. Khi đó, hiệu quả của cơ chế sửa tham số cần được đánh giá bằng tỷ lệ tác vụ hoàn tất, số lần gọi lại và khả năng kiểm soát lỗi, bên cạnh các độ đo tiền thực thi hiện có.

